\begin{abstract}

Standard Transformers have a fixed computational depth, fundamentally limiting their ability to generalize to tasks requiring variable-depth reasoning, such as multi-hop graph traversal or nested logic. We propose a depth-recurrent Transformer that decouples computational depth from parameter count by iteratively applying a shared-weight Transformer block in latent space---enabling the model to trade recurrence steps for deeper reasoning at inference time. Our architecture incorporates three mechanisms to make deep recurrence (20+ steps) stable: (1) a silent thinking objective that supervises only the final output, forcing genuine multi-step reasoning rather than intermediate heuristic shortcuts; (2) LayerScale initialization to protect fragile reasoning states from untrained layer noise; and (3) an identity-biased recurrence that creates a gradient highway across many steps. We evaluate on three compositional reasoning domains with decreasing inductive biases: graph reachability (strict adjacency masking), nested boolean logic (relative positioning), and unstructured relational text (where sequence position provides no structural hints). Across all tasks, we observe a clear \emph{computational frontier}---a boundary where performance transitions from chance to near-perfect as thinking steps scale with task complexity. Moreover, these tasks reveal qualitatively different generalization behaviors: precise but brittle (graph), approximate but robust (logic), and autonomous latent routing without structural hints (text). This progression illuminates how the interplay between a task-invariant recurrent reasoning core and task-specific perceptual interfaces shapes out-of-distribution (OOD) generalization, offering a mechanistic perspective on vertical chain-of-thought that complements the prevailing horizontal token-generation paradigm.

\end{abstract}
